%%%%%%%% ICML 2026 SUBMISSION -- LONG PAPER %%%%%%%%%%%%%%%%%%%%%
%% Workshop: Graph Foundation Models (GFM) @ ICML 2026, Seoul, July 10
%% Track: Long paper (up to 8 pages excluding refs/appendix)
%% Format: Anonymous / double-blind

\documentclass{article}

\usepackage{microtype}
\usepackage{graphicx}
\usepackage{subcaption}
\usepackage{booktabs}
\usepackage{hyperref}

\newcommand{\theHalgorithm}{\arabic{algorithm}}

% Anonymous blind-review version
\usepackage{icml2026}

\usepackage{amsmath}
\usepackage{amssymb}
\usepackage{mathtools}
\usepackage{amsthm}
\usepackage[capitalize,noabbrev]{cleveref}

\theoremstyle{plain}
\newtheorem{theorem}{Theorem}[section]
\newtheorem{proposition}[theorem]{Proposition}
\newtheorem{lemma}[theorem]{Lemma}
\newtheorem{corollary}[theorem]{Corollary}
\theoremstyle{definition}
\newtheorem{definition}[theorem]{Definition}
\newtheorem{assumption}[theorem]{Assumption}
\theoremstyle{remark}
\newtheorem{remark}[theorem]{Remark}

\icmltitlerunning{Tokenization Beats Propagation: A Position on Graph--LLM Retrieval Interfaces}

\begin{document}

\twocolumn[
  \icmltitle{Tokenization Beats Propagation:\\
    A Position on Graph Foundation Models for LLM Retrieval Interfaces}

  \icmlsetsymbol{equal}{*}

  \begin{icmlauthorlist}
    \icmlauthor{Anonymous Authors}{anon}
  \end{icmlauthorlist}

  \icmlaffiliation{anon}{Anonymous Institution}

  \icmlcorrespondingauthor{Anonymous Authors}{anon@example.com}

  \icmlkeywords{Graph Foundation Models, Retrieval-Augmented Generation, Long-Document QA,
    Graph Tokenization, LLM Reasoning, Atomic Propositions}

  \vskip 0.3in
]

\printAffiliationsAndNotice{}

\begin{abstract}
  We argue, in the context of single-document long-context QA, that the unit on
  which a graph-to-LLM retriever \emph{ranks} and the unit it \emph{reads} from
  benefit from being different objects, joined by a deterministic source-of
  map. Concretely, we treat each chunk's atomic Davidsonian proposition as a
  graph node; the source-of relation (atom $\to$ source chunk) as a deterministic
  bipartite map to the reader payload; and entity-mentions, scene index, time
  anchor, and edge-hints as edge-schema primitives. Our pipeline, Atom-Keyed
  Chunk Retrieval (AKCR), ranks atoms by query cosine, maps top atoms back to
  their source chunks, and feeds those chunks (not the atoms) to the reader.

  We evaluate on full-book NarrativeQA (40 books, 1{,}169 questions). Our main
  empirical claim is same-reader: at fixed retrieval budget $b\!=\!24$K chars,
  AKCR closes a substantial fraction of the gap to the same-reader long-context
  oracle, and outperforms a chunk-cosine retrieval baseline at the same budget,
  though it does not exceed the same-reader oracle when the oracle is
  reliable. With GLM-4.7 as reader, AKCR scores $0.5543$ vs.\ a $0.6450$
  full-book oracle (gap $-9.07$pp). On the GLM-5.1 reader, the same-reader
  comparison goes the other way ($+13.6$pp over the full-book oracle), but
  we treat this as a diagnostic of long-context degradation in GLM-5.1 rather
  than a retrieval-vs-context win. A scoring-rule ablation finds atom-best
  max-pooling the strongest of the rules we tested; alternatives degrade or
  tie within run-to-run noise. Two negative results constrain the position:
  in our single-book candidate pool, a learned typed-edge propagator
  (LearnRet) drives $\mathrm{ppr}_\alpha\!\to\!0.95$, effectively reducing to
  the dense seed; and feeding atoms directly to the reader plateaus near
  $0.46$ across four formats and two extractors.

  As a deployment-style observation we report that AKCR with Claude Opus~4.7
  at $b\!=\!24$K reaches $0.6721 \pm 0.007$ over three reruns, comparable to
  the GLM-4.7 full-book oracle ($0.6450$); we are explicit that this is a
  cross-reader comparison and not evidence that retrieval beats full context
  under a matched reader. Run-to-run variance on LLM-judged accuracy is
  non-trivial in this regime (an AKCR v1 GLM-5.1 four-run sample gave
  $\mathrm{SD}\!\approx\!5$pp on $n\!=\!1169$); we report multi-run means
  where available and treat sub-3pp single-run gaps as inconclusive. We do
  not claim graph propagation is universally unhelpful; we suggest that, on
  long single-document narrative QA, the choice of graph token is a useful
  design lever to study before adding propagation.
\end{abstract}

\section{Introduction}
\label{sec:intro}

The GFM call-for-papers calls out two themes that we address simultaneously:
\textbf{graph learning with LLMs} (``GFMs that interface with language models
and retrieval'') and \textbf{new GFMs} (``graph tokenization, structural and
positional encodings''). We take a narrow but concrete position at their
intersection: in single-document long-context QA, the unit a graph-to-LLM
retriever \emph{ranks on} and the unit the reader \emph{reads from} should
be different objects, joined by a deterministic source-of map.

Standard RAG \citep{karpukhin2020dpr} treats a chunk -- a contiguous span of
text -- as both retrieval unit and reader unit. Graph-augmented variants
\citep{edge2024graphrag,gutierrez2024hipporag} layer typed entity/relation
graphs on top of chunks and propagate via PageRank or learned GNNs. Both
families collapse two roles into one object: a structured ranking key, and
natural-language content the LLM composes from. Our position is that, at
least for long single-document narrative QA, those two roles benefit from
being represented as different units.

\paragraph{Why this is a Graph Foundation Models problem.}
We treat each chunk's set of atomic Davidsonian propositions
\citep{davidson1967logical} as the local nodes of a per-document graph. Each
atom carries an entity-mention list, a scene index, a time anchor, and a few
edge-hint relations (entity overlap, shared subject, scene co-occurrence,
inferred causal/temporal links). The source-of relation $\sigma:t\!\to\!c$
that maps an atom $t$ back to its source chunk $c$ is a deterministic
bipartite map between the graph layer and a language-payload layer. Most
prior graph-to-LLM systems decide on a graph topology and a propagation
operator first, and treat the choice of \emph{node} (chunk, entity, or
proposition) as a downstream tokenizer detail. We invert that ordering: we
fix a particular node choice (atomic propositions), make the bipartite
$\sigma$ explicit, and ask whether a learned propagator over the typed atom
graph adds anything once the dense atom seed already concentrates probability
mass on the answer region. This recasts the design question for graph-to-LLM
interfaces as ``what should the graph token be?'' before ``which propagator
do we use?''

We support this position with full-book NarrativeQA
\citep{kovcisky2018narrativeqa} evidence (40 books, 1{,}169 questions). Our
pipeline, AKCR (Atom-Keyed Chunk Retrieval), extracts $89{,}575$ atomic
propositions over $4{,}929$ chunks via a 3-pass GLM-4.7 recipe and ranks
atoms against the query in the \textsc{bge-m3} embedding space
\citep{xiao2023bgem3}. Top atoms are mapped back to source chunks via
$\sigma$ and the chunks are fed to the reader. Our main empirical claim is
same-reader (\S\ref{sec:samereader}). With GLM-4.7 as both retrieval reader
and oracle reader, HACS-AKCR at $b\!=\!24$K chars scores $0.5543$ versus a
$0.6450$ full-book oracle (gap of $-9.07$pp). At the same budget, AKCR is
above a DOS+chunks chunk-cosine baseline ($0.5227$) under the same reader,
suggesting that the gap closed against the full-context oracle is not
explained by the chunk-cosine baseline alone. With GLM-5.1 the same-reader
comparison flips ($0.5672$ AKCR vs $0.4311$ oracle on the full $\sim$78K
book); we treat that flip as a diagnostic of long-context degradation in
GLM-5.1, not as evidence that retrieval beats full context.

A scoring-rule ablation (HACS, MAR-sum, MAR-lse, ACRP) suggests atom-best
max-pooling is the strongest reduction among those we tried; alternatives
degrade or tie within run-to-run noise. Two negative results constrain the
position: a fully learned typed-edge propagator (LearnRet, APPNP closed-form
\citep{klein2024appnp}, 16 trainable parameters) drives
$\mathrm{ppr}_\alpha\!\to\!0.95$ in our single-book candidate pool, reducing
the propagation update to the dense seed; and feeding atoms themselves to
the reader plateaus near $0.46$ across four formats and two extractors.

As a separate, deployment-style observation we note that with Claude Opus
4.7 as the retrieval reader, HACS-AKCR at $b\!=\!24$K reaches
$0.6721 \pm 0.007$ over three reruns, comparable to the GLM-4.7 full-book
oracle ($0.6450$). We treat this only as a cross-reader cost-quality data
point: a stronger reader on $\sim$30\% of the book reaching the same
judge-acc as a weaker reader on the full book. It is not evidence about
retrieval vs.\ full context under a matched reader.

\paragraph{Contributions.}
We (i)~recast AKCR as a graph-foundation tokenization choice rather than a
RAG pipeline, with explicit graph-node, edge-schema, and bipartite-payload
descriptions; (ii)~report same-reader retrieval quality and oracle gaps on
NarrativeQA, where AKCR closes a substantial fraction of the gap to the
same-reader long-context oracle while not exceeding it when the oracle is
reliable; (iii)~give a scoring-rule ablation suggesting atom-best
max-pooling is the strongest reduction in this setting; (iv)~report two
negative results (learned propagation self-disable, atoms-as-content
plateau) that constrain when atomic tokenization helps and when graph
propagation adds value; and (v)~document non-trivial API run-to-run variance
for LLM-judged QA accuracy ($\mathrm{SD}\!\approx\!5$pp on $n=1169$,
GLM-5.1 reader), with a recommendation to default to multi-run reporting
and treat sub-3pp single-run gaps as inconclusive.

\section{The Position}
\label{sec:position}

\paragraph{Thesis (bounded).} \emph{For long single-document narrative QA,
our evidence suggests that the choice of graph token is a useful design
choice to study before adding propagation. In our setting, using atomic
Davidsonian propositions as ranking keys and source chunks as reader
payloads, joined by a deterministic source-of map, gives a better
budget--accuracy trade than chunk-cosine retrieval at the same budget, and
recovers a substantial fraction of full-context oracle accuracy without
exceeding it under a matched, well-behaved reader. We do not generalise to
multi-hop encyclopedic QA, schema-grounded KGQA, or non-narrative long
documents.}

\paragraph{Why role-decoupling can help here.} A long document of $C$ chunks
admits two retrieval geometries: one $d$-dimensional vector per chunk
(chunk-RAG), or roughly $17$ vectors per chunk, one per atomic proposition
(atom-RAG). Under the former, the chunk's score is the \emph{average} of
its propositions: one relevant fact buried in twenty unrelated ones is ranked
at the average. Under the latter, the chunk's score is the \emph{best} of its
propositions, because we rank atoms and only afterwards collapse to source
chunks via $\sigma$. The geometry is reminiscent of ColBERT-v2 multi-vector
retrieval \citep{khattab2022colbertv2}, with the difference that our
multi-vectors are LLM-emitted Davidsonian propositions (typed by predicate,
anchored by time and scene), not transformer token vectors.

\paragraph{Why both endpoints can fail.} A chunk is a coarse graph token: it
averages multiple propositions and discourse roles. A flat atom list is a
coarse reader payload: decontextualization \citep{choi2021decontext} strips
the narrative voice the reader composes from. We treat the atom as the
graph token, the source chunk as the language payload, and $\sigma$ as the
deterministic bipartite map between them. We do not argue this is the right
decomposition for graph-to-LLM interfaces in general; we argue it is a
practical interface choice in the single-document long-context regime, and
one whose failure cases (atoms-as-content; learned propagation on a
saturated dense seed) are themselves informative about where graph-side
operations might add value.

\section{Empirical Evidence}
\label{sec:empirical}

\subsection{Setup}
We evaluate on full-book NarrativeQA (40 books, 4{,}929 GLM-4.7-segmented chunks
of $\sim$800 chars, 1{,}169 questions). Atomic memory is extracted once, query-
independently, by a 3-pass GLM-4.7 pipeline (DEC-PROP):
(1)~\textbf{Decontextualize}, replacing pronouns with antecedent surface forms;
(2)~\textbf{Davidsonian extract}, emitting atomic propositions with explicit
\texttt{subject}, \texttt{predicate}, \texttt{object}, \texttt{time\_anchor},
\texttt{scene\_index}, \texttt{salience}, and \texttt{edge\_hints};
(3)~\textbf{Faithfulness + coverage critic}, dropping hallucinations and adding
missing atoms. Output: $89{,}575$ atoms (mean $\sim$$17$/chunk),
$\sim$$1.5\times$ the original chunk char-count, $\sim$24 minutes at 200-way
concurrency, one-time cost $\sim$\textyen 150.

For each query, AKCR scores atoms against the query using a frozen
\textsc{bge-m3} encoder, restricted to the source story; takes top-$K$ atoms
($K=200$); maps each atom to its source chunk; deduplicates while preserving
atom rank; packs source chunks DOS-sorted by chunk index into the reader's
character budget. Three readers are evaluated: Claude Opus 4.7, GLM-5.1, and
GLM-4.7. The judge is GLM-4.7 in pairwise reference-grounded mode.

\subsection{Same-reader retrieval quality and oracle gap}
\label{sec:headline}\label{sec:samereader}

We organise the empirical evidence into three blocks: (1)~same-reader
comparisons against the full-context oracle and against a chunk-cosine
retrieval baseline at matched budget; (2)~a budget sweep against the chunk
baseline under the GLM-5.1 reader; and (3)~a cross-reader deployment-style
note. Table~\ref{tab:headline} groups the same-reader rows
(\textsc{Block~A}) and a chunk-cosine baseline (\textsc{Block~B}); the
cross-reader observation is summarised in
\S\ref{sec:crossreader-note} and shown as a separate block
(\textsc{Block~C}).

\paragraph{Block A: same-reader retrieval vs.\ full-context oracle.}
At $b\!=\!24{,}000$ chars, $n\!=\!1{,}169$, with GLM-4.7 as both retrieval
reader and oracle reader, HACS-AKCR scores $0.5543$ versus a $0.6450$
full-book oracle: a gap of $-9.07$pp. On the $n\!=\!149$ subset where we
also have an Opus 4.7 oracle, AKCR scores $0.7338$ versus $0.8322$, a gap
of $-9.84$pp. We do not claim AKCR closes this gap; it does not. AKCR
recovers a fraction of full-context performance (about $86\%$ of GLM-4.7
oracle, about $88\%$ of Opus oracle on that subset) at $\sim$30\% of the
input cost. With the GLM-5.1 reader, the same-reader comparison reverses
($0.5672$ AKCR at 24K vs $0.4311$ oracle at 78K), but we treat this as
diagnostic of long-context degradation in GLM-5.1 (consistent with
lost-in-the-middle behaviour, \citealt{liu2024lostinmiddle}), not as
evidence that AKCR exceeds full context under a well-behaved reader.

\paragraph{Block B: same-reader retrieval vs.\ chunk-cosine retrieval.}
At the same $b\!=\!24$K budget, the chunk-cosine retrieval baseline
(DOS+chunks: rank chunks by $\cos(q,c)$, pack into the budget) under
GLM-4.7 scores $0.5227$, below the AKCR same-reader score of $0.5543$
(within run-to-run variance characterised in \S\ref{sec:variance}). Under
GLM-5.1, an AKCR v1 multi-run sample (four reruns, mean
$0.5552\!\pm\!0.053$) is consistent with, but not strictly above, the
single-run DOS+chunks number ($0.5389$). We therefore claim a same-reader
retrieval improvement at GLM-4.7, and a same-reader retrieval result that is
above or at parity with chunk-cosine at GLM-5.1 within noise; we do not
overclaim a robust same-reader win at GLM-5.1.

\begin{table}[t]
  \caption{Empirical results on full-book NarrativeQA, $b=24{,}000$ chars,
    $n=1{,}169$ unless noted. \textsc{Block~A} reports same-reader retrieval
    vs.\ full-context oracle; \textsc{Block~B} reports same-reader AKCR vs.\
    chunk-cosine retrieval at the same budget; \textsc{Block~C} is the
    cross-reader deployment note (\S\ref{sec:crossreader-note}). Same-reader
    AKCR remains below the same-reader oracle when the oracle is reliable
    (GLM-4.7, Opus); the GLM-5.1 oracle row is marked as diagnostic since
    it is degraded by long-context behaviour at 78K chars. ``$\pm$'' is the
    standard deviation across reruns where reported.}
  \label{tab:headline}
  \vskip 0.05in
  \centering
  \small
  \begin{tabular}{@{}llcr@{}}
    \toprule
    Method & Reader & judge\_acc & gap to same-reader oracle \\
    \midrule
    \multicolumn{4}{l}{\textsc{Block A: same-reader retrieval vs.\ full-context oracle}} \\
    Long-context oracle (78K)        & GLM-4.7 & $0.6450$           & --- \\
    HACS-AKCR ($\alpha\!=\!0.8$, 24K) & GLM-4.7 & $0.5543$           & $-9.07$pp \\
    Long-context oracle (78K, $n\!=\!149$ subset) & Opus 4.7 & $0.8322$ & --- \\
    HACS-AKCR (24K, $n\!=\!149$ subset)            & Opus 4.7 & $0.7338$ & $-9.84$pp \\
    Long-context oracle (78K) \emph{(diagnostic)}   & GLM-5.1 & $0.4311$           & --- \\
    AKCR v1 (24K, 4-rerun mean)       & GLM-5.1 & $0.5552 \pm 0.053$ & $+12.4$pp\,$\dagger$ \\
    \midrule
    \multicolumn{4}{l}{\textsc{Block B: same-reader retrieval vs.\ chunk-cosine baseline}} \\
    DOS+chunks (chunk-cosine, 24K)    & GLM-4.7 & $0.5227$           & --- \\
    HACS-AKCR ($\alpha\!=\!0.8$, 24K) & GLM-4.7 & $0.5543$           & --- \\
    DOS+chunks (24K)                  & GLM-5.1 & $0.5389$           & --- \\
    AKCR v1 (24K, 4-rerun mean)       & GLM-5.1 & $0.5552 \pm 0.053$ & --- \\
    Atoms-as-content (GAR V013, 24K)  & GLM-5.1 & $0.4594$           & --- \\
    \midrule
    \multicolumn{4}{l}{\textsc{Block C: cross-reader deployment observation}} \\
    HACS-AKCR (24K, 3 reruns)          & Opus 4.7 & $0.6721 \pm 0.007$ & --- \\
    Long-context oracle (78K)          & GLM-4.7  & $0.6450$           & --- \\
    \bottomrule
  \end{tabular}

  \medskip
  \footnotesize
  $\dagger$~\emph{Diagnostic only}: we read AKCR (24K) above the GLM-5.1
  full-context oracle (78K) as a reflection of long-context degradation in
  GLM-5.1, not as a retrieval-vs-context win.
\end{table}

\paragraph{Block~C: cross-reader deployment observation.}
\label{sec:crossreader-note}
With Claude Opus 4.7 as retrieval reader, HACS-AKCR at $b\!=\!24$K reaches
$0.6721 \pm 0.007$ over three reruns. The GLM-4.7 long-context oracle on
the full $\sim$78K-char book scores $0.6450$. We report this only as a
deployment-style data point: a stronger reader on $\sim$30\% of the book
producing comparable judge accuracy to a weaker reader on the full book.
It is not evidence about retrieval vs.\ full context under a matched reader,
and we do not treat it as a primary scientific claim. We further caution
that this observation depends on the relative reading capability of Opus
4.7 vs.\ GLM-4.7 on this benchmark; with a different reader pair, the
direction can plausibly invert.

\paragraph{Pareto across budgets.} On a single-run budget sweep
($b\!\in\!\{2,4,8,12,16,24\}$K chars, GLM-5.1 reader), AKCR v1 was above
DOS+chunks at five of six budgets, with the largest gap $+8.3$pp at
$b\!=\!12$K and a tie at $b\!=\!8$K. The associated single-run
paired-bootstrap $P$-values look strong, but because the underlying
run-to-run dispersion (\S\ref{sec:variance}) is comparable to a few of these
gaps we treat the qualitative ordering, not the per-budget significance,
as the take-away.

\paragraph{Cost note.} DEC-PROP extraction is a one-time corpus-side cost
($\sim$\textyen 150 for $4{,}929$ chunks at our concurrency setting);
\textsc{bge-m3} encoding runs on Apple MPS; per-query retrieval is a CPU
dense matmul.

\subsection{Scoring-rule ablation}
\label{sec:ablation}
The atom-keyed retrieval admits a family of chunk-scoring rules. Define the
top-$K$ atom set $T_q$ for query $q$ and let $s(t)\!=\!\cos(q,t)$ be the dense
atom-query score. We compare:
(i)~\textbf{v1}, atom-best: rank chunks by their highest-scoring atom in $T_q$;
(ii)~\textbf{MAR-sum}: chunk score $=\!\sum_{t \in c, t \in T_q} s(t)$;
(iii)~\textbf{MAR-lse} ($\tau\!=\!5$): log-sum-exp of in-chunk top-$K$ atom scores;
(iv)~\textbf{HACS} ($\alpha$): $\alpha \cdot \max_{t\in c}s(t) + (1\!-\!\alpha)\cdot\cos(q,c)$;
(v)~\textbf{ACRP}: same chunks as v1, but the top-15 atoms are also prepended
to the reader prompt as ``key facts.''
Table~\ref{tab:ablation} reports paired-bootstrap deltas vs a fresh v1 rerun
($n\!=\!1{,}169$, GLM-5.1 reader, $b\!=\!24{,}000$).

\begin{table}[t]
  \caption{Scoring-rule ablation: paired bootstrap vs v1 rerun~\#2 (May~8;
    $n\!=\!1{,}169$, GLM-5.1 reader). HACS $\alpha\!=\!0.8$ ties v1 within
    noise; alternatives degrade or tie. ACRP regresses ($-7.3$pp) because
    atoms-as-pointer pulls the reader toward atoms-as-content failure.}
  \label{tab:ablation}
  \vskip 0.05in
  \centering
  \small
  \begin{tabular}{@{}lcccr@{}}
    \toprule
    Rule & judge\_acc & $\Delta$pp & 95\% CI & $P(>v1)$ \\
    \midrule
    v1 ($\alpha\!=\!1.0$) & $0.5595$ & --- & --- & --- \\
    HACS $\alpha\!=\!0.8$ & $0.5672$ & $+0.77$ & $[-1.71,+3.25]$ & $0.72$ \\
    HACS $\alpha\!=\!0.7$ & $0.5423$ & $-1.71$ & $[-4.28,+0.77]$ & $0.08$ \\
    MAR-sum & $0.5458$ & $-1.37$ & $[-3.93,+1.11]$ & $0.13$ \\
    MAR-lse $\tau\!=\!5$ & $0.5355$ & $-2.40$ & $[-4.96,+0.17]$ & $0.03$ \\
    ACRP $K_{\mathrm{kf}}\!=\!15$ & $0.4867$ & $-7.27$ & $[-10.0,-4.62]$ & $0.000$ \\
    \bottomrule
  \end{tabular}
\end{table}

\paragraph{Reading the ablation.} In our setting, HACS $\alpha\!=\!0.8$ ties v1
(within run-to-run noise; \S\ref{sec:variance}); HACS $\alpha\!=\!0.7$, MAR-sum,
and MAR-lse all degrade weakly. ACRP regresses sharply: prepending atoms to
the reader prompt appears to induce an atoms-as-content failure mode (the
reader tends to paraphrase the atom list rather than synthesise from the
chunks). Among the rules we tried, atom-best max-pooling is the only
reduction that consistently survives. We do not claim it is uniquely optimal
in general; it is the rule we keep recommending in this single-document setup
until counter-evidence appears. HACS $\alpha\!=\!0.8$ is also the variant
that drives the cross-reader headline in Table~\ref{tab:headline}: the
$20\%$ chunk-cosine bonus appears to recover scene-/discourse-level
context without breaking the atom-dominated ranking.

\subsection{Robustness and variance-aware interpretation}
\label{sec:variance}

Across all numbers in Table~\ref{tab:headline} we apply the following
interpretation rules.

\textbf{(1)~Run-to-run variance.} Four identical-configuration AKCR v1 reruns
($n\!=\!1{,}169$, $b\!=\!24$K, \texttt{temperature=0}, GLM-5.1 reader) gave
judge\_acc $\in\!\{0.5997,0.5595,0.5825,0.4790\}$, mean
$0.5552\!\pm\!0.053$, range $0.121$. Single-run paired-bootstrap confidence
intervals are tighter than this dispersion, so we report multi-run means
where feasible and treat sub-3pp single-run gaps as inconclusive.

\textbf{(2)~Comparison to chunk baseline across budgets.} On the GLM-5.1
budget sweep, AKCR v1 was above DOS+chunks at five of six budgets. Several
of those gaps are within run-to-run noise; we therefore make a qualitative
claim (AKCR is at least at parity with chunk-cosine across budgets), not a
strong significance claim per budget.

\textbf{(3)~Oracle-gap retention.} Defining
$\mathrm{retained\_oracle\_fraction} =
\mathrm{AKCR_{24K}} / \mathrm{Oracle_{78K}}$ under a matched reader, we
obtain $\sim$$0.5543/0.6450\!\approx\!0.86$ for GLM-4.7 and
$\sim$$0.7338/0.8322\!\approx\!0.88$ for Opus on the $n\!=\!149$ subset,
i.e., AKCR retains roughly $86\%$--$88\%$ of full-context accuracy at
$\sim$$30\%$ of input chars. Defining
$\mathrm{gap\_closed\_vs\_chunk} = (\mathrm{AKCR_{24K}}\!-\!\mathrm{Chunk_{24K}})
/ (\mathrm{Oracle_{78K}}\!-\!\mathrm{Chunk_{24K}})$, we obtain at GLM-4.7
roughly $(0.5543\!-\!0.5227)/(0.6450\!-\!0.5227)\!\approx\!0.26$, i.e.,
AKCR at the matched 24K budget closes about a quarter of the gap between
chunk-cosine retrieval and the full-context oracle. Both ratios are
budget- and reader-dependent and we report them as descriptive metrics, not
as headline numbers. (\textsc{TODO}: same diagnostic at additional budgets
with paired reruns; not yet locked.)

\textbf{(4)~GLM-5.1 oracle is diagnostic.} We mark the GLM-5.1 full-context
oracle row in Table~\ref{tab:headline} as diagnostic: at 78K-char inputs
the GLM-5.1 oracle is degraded relative to GLM-4.7 in a manner consistent
with long-context behaviour reported in the literature
\citep{liu2024lostinmiddle}. Same-reader AKCR exceeding the GLM-5.1 oracle
should not be read as a general retrieval-vs-context result.

\textbf{(5)~Negative effect sizes.} The two negatives in
\S\ref{sec:negatives} have effect sizes ($|\Delta|\!\geq\!2.4$pp on the
relevant comparisons) larger than the per-run dispersion above, and the
LearnRet $\mathrm{ppr}_\alpha\!\to\!0.95$ result is structural rather than
a bootstrap effect.

\section{Lessons from Negatives}
\label{sec:negatives}

We report two negative results that together bound the position. Both are
negative outcomes for design choices a graph-foundation pipeline might try
first; we read them as constraints on \emph{when} graph-side propagation or
graph-side payloads add value, not as evidence against propagation in
general.

\subsection{Learned propagation self-disables on saturated single-book seeds}
We built a typed multi-edge graph over atoms (entity overlap, shared subject,
scene index, and edge-hint relations) and ran Personalized PageRank with
hand-tuned weights. Hand-tuned PPR lost $3.7$pp F1@15 vs.\ dense-only on the
same retrieval target. We then made all 16 graph hyperparameters trainable
as \texttt{nn.Parameter}, used the APPNP closed-form solve
\citep{klein2024appnp}, and trained with a multi-positive InfoNCE loss
(LearnRet). The optimizer drove $\mathrm{ppr}_\alpha\!\to\!0.95\!\pm\!0.01$,
which makes the propagation update reduce, in the limit, to the dense seed.
A natural interpretation is that the dense atom seed is already sufficient
when the candidate pool is restricted to a single book of $\sim$$5{,}000$
chunks: the structured graph adds little that the seed has not already
captured. We do not claim propagation is useless in general; multi-hop
encyclopedic QA, multi-document settings, and schema-grounded KGQA may
differ. We claim that, in this single-book regime with this edge typology,
the learned propagator did not find a non-trivial use of the graph.

\subsection{Atoms-as-content caps at $0.46$ regardless of format}
We tested four atom-feed-to-reader variants:
flat bullet list of V010 atoms, $0.307$ at GLM-4.7;
GAR V010 (graph format with edges), $0.458$ at GLM-5.1;
GAR V012 (richer extraction with verbatim sources and typed edges), $0.458$;
GAR V013 (verbatim source as displayed text plus modality annotations),
$0.459$. None breaks the $0.46$ ceiling. Stronger extractors do not change the
ceiling: V010 (GLM-4.7 extractor, mean lexical recall $0.859$), V011 (GLM-5.1
extractor, $0.895$), V012 (GLM-5.1 with verbatim+edges+modality, more
selective). The bottleneck is not extraction quality; it is that decontextualized
propositions lose the narrative voice the reader needs to compose answers.
Hybrid retrieval (atoms + chunks + summaries) at $b=24{,}000$ scores $0.519$;
adding RAPTOR cluster summaries does not exceed $0.431$ at $b=16{,}000$. None
of these break $0.6$.

\section{Discussion}
\label{sec:discussion}

\paragraph{One-paragraph summary.} In our setup, once the graph token is
fine-grained enough that a single dense atom seed concentrates probability
mass on the answer region, a learned propagator over typed atom edges did
not find a non-trivial use of the graph, and feeding atoms directly to the
reader hurt rather than helped. We read this as evidence that, in this
regime, the choice of graph token is worth studying before adding a
propagation operator -- not as a claim that propagation is unhelpful in
general.

\paragraph{Where this position is bounded.} Our results are restricted to
long-document QA over narrative and semi-narrative text in NarrativeQA. We
do not claim the position generalises to multi-hop encyclopedic QA
(e.g., HotpotQA, MuSiQue), where edge traversal can connect facts no single
chunk contains; to schema-grounded KGQA; or to non-narrative long documents
(legal, scientific, code). Each is a candidate stress test and a candidate
falsifier. In particular, distractor-heavy multi-hop benchmarks such as
HotpotQA are a natural future stress test for AKCR: when distractor
paragraphs are crafted to share entities with the query, a max-pool over
atom cosines is no longer obviously the right ranking signal, since
high-similarity atoms may come from distractors. We frame that as a
hypothesis to be tested, not as a reported result of this paper.

\paragraph{Open questions for the GFM community.}
(i)~\emph{What is the right tokenization granularity, and how should it
scale with document length?} Atomic propositions work in our setup;
sub-atomic fragments may overshoot. A document-length-dependent
``atoms-per-chunk'' study, with consistent extraction quality, would be
useful.
(ii)~\emph{When does graph propagation add value on top of a strong dense
atomic seed?} Our learned propagator self-disabling is a falsification
signal, not a general refutation; multi-hop, multi-document, or
sparse-coverage regimes are the natural places to look for the boundary.

\section{Limitations}
We evaluate on a single benchmark family (NarrativeQA full-book) and three
commercial/API readers (GLM-4.7, GLM-5.1, Claude Opus 4.7); the
generalisation in \S\ref{sec:discussion} remains a hypothesis until tested
on multi-hop QA, schema-grounded KGQA, and non-narrative long documents.
Our negative result on learned propagation is specific to a single-book
candidate pool, this particular edge typology, and our APPNP-style
parameterisation; a different graph construction or a multi-document pool
may reach a different conclusion. The DEC-PROP extractor is itself an LLM
and inherits LLM extraction artifacts; we mitigate with a
coverage-and-faithfulness critic pass but cannot guarantee zero hallucinated
or missed atoms. The judge is GLM-4.7; we use a shared judge across all
conditions to control for judge bias, but the absolute number $0.6450$ is
judge-conditional and should not be compared directly to NarrativeQA scores
in prior work using a different judge. Finally, our reader pool is
predominantly commercial-API; a similar evaluation with open-weight readers
(e.g., Llama or Qwen-class models) and a wider variance budget is a
necessary check we have not yet done.

\section*{Declaration on LLM Use}
LLMs are used as pipeline components: \textbf{GLM-4.7} for DEC-PROP atomic
extraction (3 passes), as one of three reader configurations, and as the
shared pairwise reference-grounded judge across all conditions; \textbf{GLM-5.1}
as a second reader configuration and as an alternative extractor in
ablation variants (not as judge); \textbf{Claude Opus 4.7} as the third
reader configuration. All three are commercial/API models accessed during
the experimental period. The authors used automated tools for copy-editing
and LaTeX formatting assistance. All experiments, numerical results,
tables, claims, and interpretations were verified by the authors and
remain their responsibility.

% \section*{Acknowledgements}
% Acknowledgements omitted for double-blind review.

% Ensure all references are emitted in the bibliography even when not
% cited inline. Keep citation-key list in sync with references.bib.
\nocite{khattab2022colbertv2,choi2021decontext,davidson1967logical,
  edge2024graphrag,gutierrez2024hipporag,sarthi2024raptor,klein2024appnp,
  karpukhin2020dpr,kovcisky2018narrativeqa,liu2024lostinmiddle,
  xiao2023bgem3,reimers2019sentencebert}

\bibliographystyle{icml2026}
\bibliography{references}

\end{document}
